\begin{figure}
    \centering
    \includegraphics[width=0.99\textwidth]{content/figures/hdcells.pdf}
    \caption{\textbf{Individual differences in head-direction coding around a shared ring manifold}. \textbf{a.} Left, example neurons. Right, all neurons sorted (cross-validated across time) by their preferred direction angle. Bottom, top three principal components of all neurons reveal a high-dimensional ring. \textbf{b.} Left, illustration of the cross-validation approach used for this dataset. Right, independent folds of animal A (blue and purple, top) have smaller mean squared-error (MSE) distance after alignment than a similar analyses between animal A and B (blue and red, below). \textbf{c.} Same analyses in b), repeated for all animals and 1 individual split. Distance between two different splits of the same animal are significantly smaller than across different animals (t=-3.13, p=0.0017). Middle, the former replicated on n=100 different splits. Whitening the data before alignment captures individual differences slightly better (p=0.05). MDS of the dissimilarity matrix suggests that tuning width drives across-animals differences (red gradient vs green gradient). Inset, decoding $R^2$ (using bi-cross validation in b) of tuning width from this space is significantly above chance (95\% C.I. of shuffle decoder in gray). \textbf{d.} Analyses on a separated dataset of n=6 animals shows that a rotation in the cue space (90\textdegree) is reflected in a consistent angle rotation in the neural space (98\textdegree). \textbf{e.} Left, aligning all subjects (n=31) to a common reference (Methods) reveals a high-dimensional ring. Bottom, the first principal components of the aligned space (dark blue) explains a substantial amount of variance of the original space (light blue). Right, illustration of a 0-shot learning decoding procedure form 2 aligned animals. As we aligned a larger number of subjects, the decoder approaches the performance of the decoder trained and tested on independent folds of the same animal. }
    \label{fig:hdcells}
\end{figure}

\subsection{Animal-to-animal variability in head direction coding}

Having found that population geometry is organized across brain regions, we now turn to a another scale: variability across individuals. Specifically, we ask whether the geometry of a single region differs systematically from one animal to the next. With this aim, we use shape distances to illuminate animal-to-animal variability in two recently published datasets of head direction cells (\cite{duszkiewicz}).
In both datasets, spiking activity in the mouse postsubiculum (PoSub) was collected using silicone probes during open-field exploration while position and heading direction were captured by an overhead video camera (for details, see \cite{duszkiewicz}).
The first dataset consisted of $K = 31$ sessions / subjects and a total of $2691$ neurons.
We first collected single-unit tuning curves for heading across all sessions and sorted them according to peak activity.
Preferred heading directions were roughly uniformly distributed (\textit{Fig.} \ref{fig:hdcells}a) and, applying \textit{PCA} to this aggregated dataset revealed a high-dimensional ring structure (\textit{Fig.} \ref{fig:hdcells}a, bottom), consistent with the attractor manifold previously described in this circuit \parencite{chaudhuri2019intrinsic} and akin to the simulated dataset we analyzed above.

To our knowledge, it is unknown whether there are systematic animal-to-animal differences in the statistics of head direction tuning curves in PoSub.
To investigate this, we measured the Procrustes distance between two disjoint partitions of the data along neuron and time dimensions (see illustration in \textit{Fig.} \ref{fig:hdcells}b, left; \textit{Methods}).
On average, the distance between partitions of the same animal was smaller than the distance between partitions taken from different animals (\textit{Fig.} \ref{fig:hdcells}b-c).
Whitening the data before alignment, slightly improved the capture of individual differences (\textit{Fig.} \ref{fig:hdcells}c, middle). 

To investigate what features of the neural tuning curves account for these individual differences in ring shape, we first used unsupervised methods to perform an exploratory analysis.
Akin to our analysis of the simulated data above (\textit{Fig.}~\ref{fig:synth}b-c), we used MDS followed by PCA to visualize each subject within a 2D space.
MDS introduced a median distortion \textapprox{}2\% and PCA captured \textapprox{64}\% of variance in the MDS embedding.
Coloring each session by the average width and average amplitude of the head direction tuning curves suggested that the former, and not the latter, was a principal driver of inter-animal differences (\textit{Fig.} \ref{fig:hdcells}c, right). 
To confirm this finding in our exploratory analysis, we turned to supervised methods.
The decoding of the tuning width from this space was significantly above zero (bi-cross validated as in \textit{Fig.} \ref{fig:hdcells}b), demonstrating that indeed tuning width is driving the inter-individual differences captured by the shape metrics (\textit{Fig.} \ref{fig:hdcells}c, inset). We found similar signatures of individual differences in a second, smaller dataset ($K = 6$ animals).
In this experiment, open exploration was performed in boxes with a salient landmark which was rotated by 90 degrees during the course of the recording (cue 1 vs cue 2, \textit{Fig.} \ref{fig:hdcells}d).
We find that the optimal rotation identified by Procrustes analysis in neural firing rate space was roughly matched in magnitude to the cue rotation (\textit{Fig.} \ref{fig:hdcells}d).

The analyses above show how the Procruste distance can be used to detect individual differences.
We now show that a similar analysis can identify common manifold structure across subjects and leverage this to decode head direction in one heldout subject from training data recorded in other subjects.
We started by using rotation and reflections to align all subjects to a common reference---a procedure known as \textit{generalized Procrustes analysis} \parencite{gower1975generalized}.
This revealed a high-dimensional ring structure (\textit{Fig.} \ref{fig:hdcells}e), similar to performing PCA on all data (\textit{Fig.} \ref{fig:hdcells}a) but now confirming that this structure is shared by all subjects and not just present in a subset of them. In fact, the first principal components of this aligned space explained significantly less amount of variance than the original space of each animal (\textit{Fig.} \ref{fig:hdcells}e, left), suggesting individual differences despite a common shape across subjects. Because all animals are aligned to a common shape (in this case, a ring), we could train a decoder on a subset of animals and predict head direction angles on left-out subjects, essentially performing zero-shot learning (\textit{Fig.} \ref{fig:hdcells}e, right). An analogous strategy has been used to decode movement across 4 individuals from aligned latent dynamics in motor cortex \parencite{safaie2023preserved}.

In sum, all subjects encode the head direction angle in a ring, but each of these rings have idiosyncratic distortions in higher dimensions. Both common structures and individual differences can be analyzed using shape metrics. 

